\begin{flushright}
  \fechaclase{14 de septiembre de 2026}
\end{flushright}

\subsection{Seguimiento de trayectorias por SMC}

Considere el sistema de segundo orden
\[
  \dot{X}_1=X_2
\]
\[
  \dot{X}_2=f(X,t)+v
\]
\[
  y=X_1
\]
donde la dinámica no lineal cumple
\[
  \lvert f(X,t)\rvert\leq L
\]
para algún real positivo $L$.

\paragraph{Definición}
Para el sistema SISO con salida $y\in\mathbb{R}^{n}$ y entrada
$v\in\mathbb{R}$. Si la $i$-ésima derivada de la salida $y^{(i)}$ es
independiente de la entrada $v$ para toda $i=1,2,\ldots,k-1$, pero $y^{(k)}$
es proporcional a $v$ en un dominio $\Omega\in X$, se dice que $k$ es el
grado relativo bien definido del sistema.

El problema de control consiste en hacer que la salida $y$ siga una señal de
referencia $r$. El sistema es de grado relativo $k=2$, pues
\[
  y'=\dot{X}_1=X_2
\]
\[
  y''=\dot{X}_2=f(X,t)+v
\]
el error de seguimiento se define como
\[
  e=r-y=r-X_1
\]

Para hacer que $y\to r$ se busca que $e\to0$ mediante una función de
Lyapunov, por lo que se propone una superficie de deslizamiento como una
ecuación diferencial estable de grado $k-1$, entonces
\[
  \sigma=\dot{e}+me
\]
y derivando los errores
\[
  \dot{e}=\dot{r}-\dot{y}=\dot{r}-X_2
\]
\[
  \ddot{e}=\ddot{r}-\dot{X}_2
  =\ddot{r}-f(X,t)-v
\]

Se propone la función de Lyapunov
\[
  V=\frac{1}{2}\sigma^2
\]
y se deriva la superficie
\begin{align*}
  \dot{\sigma}
  &=m(\dot{r}-X_2)+\ddot{r}-f(X,t)-v\\
  &=m\dot{r}+\ddot{r}-mX_2-f(X,t)-v
\end{align*}

Considerando que las derivadas de la referencia son acotadas
\[
  \lvert m\dot{r}+\ddot{r}\rvert\leq R
\]
para algún real positivo $R$, de modo que
\[
  \dot{V}=\sigma\dot{\sigma}
  =\sigma\left[m\dot{r}+\ddot{r}-f(X,t)-mX_2-v\right]
\]
\[
  \dot{V}\leq\sigma\left[R+L-mX_2-v\right]
\]

Se elige
\[
  v=-mX_2+K\operatorname{sign}(\sigma)
\]
\[
  \dot{V}\leq\sigma\left[R+L-K\operatorname{sign}(\sigma)\right]
\]
\[
  \dot{V}\leq(R+L)\lvert\sigma\rvert-K\lvert\sigma\rvert
\]
\[
  \dot{V}\leq-(K-R-L)\lvert\sigma\rvert
\]

Si $K>R+L$ el error de seguimiento es asintóticamente estable y $e\to0$
conforme $t\to\infty$ con la ley de control
\[
  \sigma=\dot{e}+me
\]
\[
  v=-mX_2+K\operatorname{sign}(\sigma)
\]